\section{Conclusion and Future Work}
\label{ch:conclusion}

\subsection{Conclusions}

We set out to catch a developing engine fault from the standard 1\,Hz
OBD-II stream, before a diagnostic trouble code fires. What we built is a
physics-grounded fault injection engine for four common fault classes; a
vehicle-agnostic feature and normalization pipeline that lets a model
trained on one car run on another; a six-class XGBoost classifier evaluated
under session-level and family-level cross-validation; a per-PID
forecaster and a one-class anomaly detector as a second, independent line
of evidence; a native Android application, OBD-LOG, for real-vehicle data
collection over Bluetooth; a live Streamlit dashboard that explains its own
predictions and closes every session with a plain-language health verdict;
and, threaded through all of it, a reporting layer that says ``Untested''
or ``insufficient data'' instead of manufacturing a number it cannot back.

What we did not build is a validated real-fault detector. That distinction
matters more than any F1 score in this book, and we have tried not to blur
it anywhere above. Every classification number in Chapter~\ref{ch:results}
measures the model's ability to recover labels written by our own
injector. Building the validated version needs one thing our project could
not supply on its own timeline: labeled recordings of real, induced faults.
Section~\ref{sec:future} makes that the first line of future work, not the
last, because it is the single change that would do the most.

If this project has one engineering contribution worth naming beyond the
models themselves, it is the discipline of refusal. Twice in this book the
system was handed a situation it could not honestly resolve --- a vehicle
missing the sensor a fault class depends on, a recording too short to
support a verdict --- and twice it said so instead of guessing. The
alternative version of this system, the one we shipped first and then
un-shipped, guessed with 97\,\% confidence on a healthy car. We consider the
version that admits what it does not know to be the more useful tool, even
though it produces fewer confident-looking numbers for a slide.

That discipline did not come free. It cost us a working demo we had to
throw away, a redesign of the reporting layer around an evaluability
contract nobody had asked for at project kickoff, and an uncomfortable
week of explaining to each other why a 97\,\% confidence score was the
project's worst result rather than its best. We record that cost here
because a reader building a similar system will hit the same choice ---
ship the number the model gives you, or build the machinery to know when
that number is not trustworthy --- and the second path is more work for a
less impressive-looking demo. We think it is the only defensible path for
anything a mechanic or a driver might actually act on.

The broader lesson carries past this particular car and these particular
four faults. Any system trained on synthetic labels inherits the honesty
of its label generator, and any system deployed across more than one
instance of its target domain inherits the assumptions baked into its
first training example. We tried to make both of those inheritances
visible rather than implicit: the injector's physics reasoning lives next
to its code, the withheld-configuration gap is a number in this book and
not a footnote, and the cross-vehicle failure that produced the Untested
contract is documented as a failure, not smoothed into a success story
after the fact. If a future team extends this work, we would rather they
inherit the visible assumptions than have to rediscover them the way we
did.

\subsection{Limitations}

We list these without softening them, because the audience for this report
already knows the honest framing and overselling would cost more credibility
than it could ever buy back.

\begin{enumerate}
  \item \textbf{Every fault metric in this book measures self-consistency,
    not field performance.} The injector writes the fault; the classifier
    learns to invert the injector. A macro-F1 of 0.958 on a held-out
    driving family says the model generalizes across driving contexts for
    faults it was told to recognize. It does not say the model would catch
    a real vacuum leak on a real car. No experiment in this project
    contradicts that limitation, and none was designed to.
  \item \textbf{The model has never detected a real fault.} Not once, on
    any vehicle, in any recording. The real-fault evaluation code exists
    in the repository and its data folder is empty. This is the single
    most important gap between what we built and what a mechanic would need
    to trust it.
  \item \textbf{Air-system detection is unavailable on the majority of
    modern cars.} MAF-based engines --- our own test vehicle among them ---
    expose no MAP sensor, and the air-system fault's primary signature
    lives on MAP. On those cars the class is correctly marked Untested, but
    correctly refusing to answer is not the same as answering.
  \item \textbf{Phone-based sampling aliases fast dynamics.} Our measured
    ELM327 poll rate falls to 0.37\,Hz at 21 polled PIDs --- well under the
    1\,Hz the training pipeline assumes. Results scored on app-captured
    data should be read as pessimistic relative to a faster logger, because
    slow polling smears short fuel-trim oscillations into apparent
    stability, which can understate real trim variance and, in the other
    direction, overstate false-positive risk on transient events the
    sampler simply missed.
  \item \textbf{Self-baselining cannot see a fault that predates the
    recording.} The end-of-read report treats the first 20\,\% of a session
    as healthy ground truth. A car that is already faulty at key-on has
    that fault baked into its own baseline and will read as normal for the
    rest of the drive. This is a structural property of the self-baseline
    design, not a bug to be patched.
\end{enumerate}

\subsection{Future Work}
\label{sec:future}

The items below are ordered by impact, not by how soon they could be
started. Each one names a specific mechanism and a specific reason it was
not already done, because ``improve the model further'' is not a plan.

\begin{enumerate}
  \item \textbf{Real fault data.} This is the one change that would convert
    the project from ``a detection algorithm plus a collection protocol''
    into a validated detector. A single matched set of before, during, and
    after recordings of a genuinely induced fault --- a vacuum line
    partially disconnected under a mechanic's supervision, say --- would let
    the real-fault evaluation code that already exists in the repository
    produce its first non-empty result. We did not do this in the current phase
    because it requires a donor vehicle, a supervising mechanic, and a
    written abort protocol we were not yet in a position to execute safely;
    the collection protocol documenting exactly how to do it is already
    written and waiting in the repository, alongside the scorer that would
    consume its output.

  \item \textbf{A live Android-to-server pipeline.} Route the app's CSV
    output over HTTPS, in roughly 60-second chunks, to a cloud-hosted
    version of the existing FastAPI backend, which already implements JWT
    authentication. The server holds one persistent scoring session per
    vehicle --- one fitted baseline, every chunk scored against it in
    sequence --- and the Streamlit dashboard mirrors that session live
    while the phone shows a simplified health card. This was not built now
    because it needs a hosting decision (which cloud, which uptime
    guarantee) that is out of scope for a laptop-based thesis defense, and
    because it depends on item 1 producing enough real sessions to be
    worth streaming.

  \item \textbf{Apache Kafka for multi-vehicle ingest.} At fleet scale,
    direct HTTP ingest does not survive a phone losing signal mid-upload,
    and a new server process per vehicle does not scale. The natural next
    architecture makes the Android app a Kafka producer, runs a consumer
    group that performs inference per vehicle, and writes a second consumer
    to a time-series store for the dashboard to query. Kafka's offset-based
    replay absorbs cellular dropouts that a plain HTTP POST cannot, and a
    Kafka Connect OBD-II source connector is the natural long-term target.
    This is future work rather than current work because it only pays for
    itself once item~2 exists and more than one or two vehicles are
    actually streaming --- building it first would be premature
    infrastructure for a single-laptop demo.

  \item \textbf{Per-regime baselining.} Right now one healthy baseline
    covers idle and driving alike, and we measured the resulting
    idle-to-drive feature shift at a mild 2--3$\sigma$, concentrated on the
    throttle side, and already absorbed by the regime one-hot features and
    the alerter's temporal vote. Because the measured shift is mild, this
    is a robustness improvement queued behind the higher-impact items
    above, not a critical bug demanding an emergency fix. Implementing it
    properly means retraining with regime-conditioned normalization rather
    than patching the existing scaler.

  \item \textbf{Intermittent and accelerating fault shapes in the
    injector.} Every fault the injector produces today follows one shape:
    a linear ramp to a held plateau. Real components do not always fail
    that way --- a connector that only opens under vibration produces an
    intermittent signature, and bearing-style wear often follows a
    Weibull-shaped acceleration rather than a straight line. Widening the
    injector's shape library would reduce the risk that the classifier is
    quietly overfitting to ``ramp-then-hold'' as a pattern rather than
    learning the underlying fault physics. We scoped this out of the
    current phase because the four existing fault models, in ramp and step
    mode, already consumed the available time for physics validation, and
    adding shapes without re-validating physics for each one would
    reintroduce exactly the kind of unaudited synthetic data this project
    was built to avoid.

  \item \textbf{A session-label field in the Android app.} A single text
    field, written as one more CSV column, so a recording can declare
    itself \emph{baseline}, \emph{fault-type-X}, or \emph{healthy-post-repair}
    without an external notebook tracking which file is which. It is a
    small UI change, deliberately left until last because it only becomes
    necessary once real-fault collection (item~1) is actually underway ---
    there is no annotation burden to solve before then.
\end{enumerate}
